\documentclass[10pt,a4paper]{amsart}
\usepackage[margin=19mm]{geometry}
\usepackage{amsmath,amssymb,mathtools}
\usepackage[scr=rsfs,cal=euler]{mathalfa}
\usepackage{xfrac}
\usepackage[unicode,hidelinks,bookmarks=false]{hyperref}
\newcommand{\ie}{i.e.\ }
\newcommand{\cf}{cf.\ }
\newcommand{\resp}{respectively\ }
\newcommand{\Subsection}{\subsection}
\providecommand{\nobreakdash}{\nobreak\hbox{-}}
\setlength{\emergencystretch}{2em}
\multlinegap0pt
\usepackage{bm}
\usepackage{bbm}
\usepackage{dsfont}
\usepackage{xspace}
\usepackage{cancel}
\usepackage{mathtools}
\usepackage{amsthm}
\makeatletter
\@ifundefined{theorem}{\newtheorem{theorem}{Theorem}}{}
\@ifundefined{lemma}{\newtheorem{lemma}[theorem]{Lemma}}{}
\@ifundefined{proposition}{\newtheorem{proposition}[theorem]{Proposition}}{}
\makeatother
\DeclareMathOperator{\ord}{ord}
\newcommand{\Dt}{\Delta}
\newcommand{\calL}{{\mathcal{L}}}
\newcommand{\calT}{{\mathcal{T}}}
\newcommand{\dd}{\delta}
\title[The Equivalence Problem for Generalized Airy Operators]{The Equivalence Problem for Generalized Airy Operators}
\author{Yichuan Cao, Ruyong Feng, Yunfei Li, Ruichen Qiu}
\date{Source: arXiv:2607.00359v1 (CC BY 4.0)}
\begin{document}
\maketitle

\noindent\emph{Source and licence.} The Equivalence Problem for Generalized Airy Operators. Version arXiv:2607.00359v1 (CC BY 4.0), distributed under the Creative Commons Attribution 4.0 International licence, as recorded in the arXiv licence metadata for this exact version and shown on the abstract page https://arxiv.org/abs/2607.00359v1. Full rights notice: licenses/arxiv-2607.00359-CC-BY-4.0.txt.\par\smallskip

\noindent\textit{Editorial scope.} Counted unit: the Lemma whose source label is \texttt{LM:minimalorderinL2} in the article (source0.tex lines 677--697, source label \texttt{LM:minimalorderinL2}), delivered together with the article's own complete original proof of it (source0.tex lines 698--831, one \texttt{proof} environment of 134 lines). No mathematics is written, repaired or re-derived: statement and proof are the article's own text line for line. Required context retained uncounted: PROP:Puiseuxsolutions (lines 183--194); LM:logderivative (lines 519--528); LM:order-polynomialdifference (lines 551--557); LM:Bellformula (lines 589--597); LM:orderofBell1 (lines 607--613); LM:orderofBell2 (lines 634--641). Every definition, display and label of those lines is retained unchanged. Omitted from this package and not redistributed: 1--182, 195--518, 529--550, 558--588, 598--606, 614--633, 642--676, 832--1059. Nothing inside a retained excerpt is omitted or altered: every retained body is delivered exactly as the article's lines are written, so no omission receipt of any kind accompanies this package. Citations: the retained excerpts carry no citation command at all, so no bibliography entry of the article is transcribed and no reference is redistributed. Reference adaptation: none was needed and none was made; every reference inside the delivered unit resolves inside it, so no unresolved reference or citation is produced. Printed slips kept literal: no printed word, symbol, hypothesis or number of the counted statement or of its proof was altered. Layout and class: the article's own class is \texttt{article} with packages including authblk, amsmath, amssymb, amsthm, thmtools; the wrapper is the lane's pinned \texttt{amsart} editorial wrapper at 10pt on A4, which restates the theorem-like environments the retained text needs and adds the article's own macro definitions that the retained text needs (\texttt{\textbackslash Dt}, \texttt{\textbackslash calL}, \texttt{\textbackslash calT}, \texttt{\textbackslash dd}, \texttt{\textbackslash ord}), with extra wrapper packages (bm, bbm, dsfont, xspace, cancel, mathtools). The delivered numbering is the unit's own. Assets: no retained span contains a figure environment, a graphics-inclusion command, a PDF-page inclusion, a table or a TikZ picture; no image file is copied into this package, the package declares no asset dependency and no image file is redistributed with it. Licence: version arXiv:2607.00359v1 of this article is distributed under the Creative Commons Attribution 4.0 International licence, as recorded in the arXiv licence metadata for this exact version and shown on the abstract page https://arxiv.org/abs/2607.00359v1. A full rights notice is delivered at licenses/arxiv-2607.00359-CC-BY-4.0.txt. Characters: every retained excerpt is pure ASCII, so the delivered engine, XeLaTeX with its default Latin Modern text font, reports no missing character.
\begin{proposition}
\label{PROP:Puiseuxsolutions}
Let $P\in k[Y]$ be monic of degree $n>0$, and let $Q\in k[x]$ have degree $m>0$ and leading coefficient $b_m$. Put
\[
  e=\frac{n}{\gcd(n,m)},\qquad q=\frac{m}{\gcd(n,m)},\qquad v=t^{1/e}.
\]
If $c_1,\dots,c_n$ are the roots of $W^n+b_m=0$, then there exist unique units $w_{c_i}(v)\in k[[v]]$ with $w_{c_i}(0)=c_i$ such that
\[
  v^{-q}w_{c_1}(v),\dots,v^{-q}w_{c_n}(v)
\]
are precisely the roots of $P(Y)+Q(x)=0$ in $\bigcup_{r=1}^\infty k((t^{1/r}))$.
\end{proposition}

\begin{lemma}
\label{LM:logderivative}
Let $w$ be a nonconstant Puiseux Laurent series.  Then 
$$
 \ord_t\!\left(\frac{\dd (w)}{w}\right)\ge 1.
$$ 
Moreover, 
$\ord_t\!\left(\frac{\dd w}{w}\right)=1$ if and only if $\ord_t(w)\neq 0$.

\end{lemma}

\begin{lemma}
  \label{LM:order-polynomialdifference}
Let $P(Y) \in k[Y]$ be a monic polynomial of degree $n$, and $w, \Dt \in k((t^{1/e}))$ be two Puiseux series. 
If $\ord_t(w) < 0$ and $\ord_t(\Dt) > \ord_t(w)$, then  
$$P(w+\Dt)-P(w) = nw^{n-1}\Dt + M,$$
where $\ord_t(M)>(n-1)\ord_t(w)+\ord_t(\Delta)$.
\end{lemma}

\begin{lemma}
  \label{LM:Bellformula}
Suppose that $\calL=a_n\partial^n+a_{n-1}\partial^{n-1}+\dots+a_0\in R$.
Then 
$$
   \frac{\calL(z)}{z}=\sum_{i=0}^n a_i B_i(\gamma)
$$
where $\gamma=\frac{\delta(z)}{z}$.
\end{lemma}

\begin{lemma}
  \label{LM:orderofBell1}
  Suppose that $w\in k((t^{1/e}))$ with $\ord_t(w)<1$. Then for $j\geq 0$
  $$
     \ord_t(C_j(w))>j\ord_t(w).
  $$
\end{lemma}

\begin{lemma}
  \label{LM:orderofBell2}
  Suppose that $w_1,w_2\in k((t^{1/e}))$. Suppose that 
  $\ord_t(w_1)<\ord_t(w_2)<\infty$ and $\ord_t(w_1)<0$. Then for $j\geq 0$
  $$
     \ord_t(C_j(w_1+w_2)-C_j(w_1))> (j-1)\ord_t(w_1)+\ord_t(w_2).
  $$
\end{lemma}

\begin{lemma}
  \label{LM:minimalorderinL2}
  Suppose that $\calL_1=P_1(\partial)+Q_1(x)$ and $\calL_2=P_2(\partial)+Q_2(x)$
  are two generalized Airy operators of type $(n,m)$ with $nm\neq 0$. Assume that 
  $z=\exp^{\int \eta dx} h$ with $\eta, h\in k((t^{1/e}))$ satisfying that
  $\calL_1(z)=0$ and $P_1(\eta)+Q_1(x)=0$. Let $\calT\in R$ satisfy $\calT(z)/z \notin k$, and set $\alpha=\frac{\delta(z)}{z}, \beta=\frac{\delta(\calT(z))}{\calT(z)}$. 
  Then the following assertions hold
  \begin{enumerate}
  \item [(1)] $\ord_t(P_2(\beta)-P_2(\alpha))=(n-1)\ord_t(\alpha)+\ord_t(\beta-\alpha).$
    \item [(2)]
    $\frac{\calL_2(\calT(z))}{\calT(z)}=(Q_2-Q_1)
 +(P_2(\alpha)-P_1(\alpha))
 +\bigl(P_2(\beta)-P_2(\alpha)\bigr)+M
  $
  \item [(3)]$\ord_t(M)>
     \min
     \{\ord_t(P_2(\alpha)-P_1(\alpha)), \ord_t(Q_2-Q_1),
     \ord_t(P_2(\beta)-P_2(\alpha))\}.
  $
  \end{enumerate}
\end{lemma}

\begin{proof}
Set $\Delta=\beta-\alpha$.
Since $\Delta=\delta(\calT(z)/z)/(\calT(z)/z)$ and $\calT(z)/z\notin k$,  Lemma~\ref{LM:logderivative} gives
$
  \ord_t(\Delta)\geq 1.
$
On the other hand,
\[
  \alpha=\eta+\frac{\delta (h)}{h}.
\]
By Proposition~\ref{PROP:Puiseuxsolutions} and 
Lemma~\ref{LM:logderivative}, $\ord_t(\eta)=-m/n<0$ and
$\ord_t(\delta(h)/h)\geq 1$. Hence we have
$
  \ord_t(\alpha)=-\frac{m}{n}<0.
$
Consequently,
$$
  \ord_t(\Delta)>\ord_t(\alpha).
$$
(1) Because $P_2$ is monic of degree $n$ and $\beta=\alpha+\Delta$, by Lemma~\ref{LM:order-polynomialdifference},
$$
   P_2(\beta)-P_2(\alpha)=P_2'(\alpha)\Delta+\text{higher order terms}.
$$
Therefore,
\[
  \ord_t(P_2(\beta)-P_2(\alpha))
  =\ord_t(P_2'(\alpha)\Delta)=(n-1)\ord_t(\alpha)+\ord_t(\Delta).
\]


(2) Write
$
  P_i(Y)=\sum_{j=0}^n p_{i,j}Y^j, i=1,2.
$
By Lemma~\ref{LM:Bellformula}, we have 
\[
  0=\frac{\calL_1(z)}{z}
   =Q_1+\sum_{j=0}^n p_{1,j}B_j(\alpha)=Q_1+P_1(\alpha)+\sum_{j=0}^n p_{1,j}C_j(\alpha)
\]
and
\[
  \frac{\calL_2(\calT(z))}{\calT(z)}
   =Q_2+\sum_{j=0}^n p_{2,j}B_j(\beta)=Q_2+P_2(\beta)+\sum_{j=0}^n p_{2,j}C_j(\beta).
\]
Subtracting the first identity from the second and using
$\beta=\alpha+\Delta$, we obtain
\[
\begin{aligned}
  \frac{\calL_2(\calT(z))}{\calT(z)}
  ={}&(Q_2-Q_1)+(P_2(\alpha)-P_1(\alpha))
     +\bigl(P_2(\beta)-P_2(\alpha)\bigr)\\
   &+\sum_{j=0}^n(p_{2,j}-p_{1,j})C_j(\alpha)
     +\sum_{j=0}^n p_{2,j}
       \bigl(C_j(\beta)-C_j(\alpha)\bigr).
\end{aligned}
\]
Thus we may take
\[
  M=M_1+M_2,
\]
where
\[
  M_1=\sum_{j=0}^n(p_{2,j}-p_{1,j})C_j(\alpha)\,\,\text{and}\,\,
  M_2=\sum_{j=0}^n p_{2,j}
       \bigl(C_j(\beta)-C_j(\alpha)\bigr).
\]
This proves (2). 


(3) By Lemma~\ref{LM:orderofBell2}, for every nonzero summand of $M_2$,
\[
\begin{aligned}
  \ord_t\bigl(C_j(\beta)-C_j(\alpha)\bigr)
  &>(j-1)\ord_t(\alpha)+\ord_t(\Delta)\geq (n-1)\ord_t(\alpha)+\ord_t(\Delta),
\end{aligned}
\]
because $j\leq n$ and $\ord_t(\alpha)<0$. Together with (1), this yields
\begin{equation}
  \label{EQ:M2}
  \ord_t(M_2)>
  \ord_t\bigl(P_2(\beta)-P_2(\alpha)\bigr)
\end{equation}

Suppose that $P_1=P_2$. Then $M_1=0$. Hence
$$
  \ord_t(M)=\ord_t(M_2)>\ord_t\bigl(P_2(\beta)-P_2(\alpha)\bigr)
$$
and the desired inequality follows. 

Suppose that $P_1\neq P_2$. Because
$\ord_t(\alpha)<0$,
$$
  \ord_t(P_2(\alpha)-P_1(\alpha))=r\ord_t(\alpha).
$$
On the other hand, by Lemma~\ref{LM:orderofBell1}, for every nonzero summand of $M_1$, 
\[
\begin{aligned}
  \ord_t\bigl((p_{2,j}-p_{1,j})C_j(\alpha)\bigr)
  >j\ord_t(\alpha)
  \geq r\ord_t(\alpha)
   =\ord_t(P_2(\alpha)-P_1(\alpha)),
\end{aligned}
\]
because
$\ord_t(\alpha)<0$ and $j\leq r$. Hence
\begin{equation}
  \label{EQ:M1}
  \ord_t(M_1)> \ord_t(P_2(\alpha)-P_1(\alpha)).
\end{equation}

From (\ref{EQ:M1}) and (\ref{EQ:M2}), and the ultrametric inequality,
\[
\begin{aligned}
  \ord_t(M)
  &\geq \min\{\ord_t(M_1),\ord_t(M_2)\}\\
  &>\min\Bigl\{
      \ord_t\bigl((P_2-P_1)(\alpha)\bigr),
      \ord_t\bigl(P_2(\beta)-P_2(\alpha)\bigr)
     \Bigr\}.
\end{aligned}
\]
Since adjoining the additional number $\ord_t(Q_2-Q_1)$ can only decrease
this minimum, we conclude that
\[
\begin{aligned}
  \ord_t(M)>
  \min\bigl\{&\ord_t\bigl((P_2-P_1)(\alpha)\bigr),
  \ord_t(Q_2-Q_1),
  \ord_t\bigl(P_2(\beta)-P_2(\alpha)\bigr)\bigr\}.
\end{aligned}
\]
This proves the case $P_1\neq P_2$. Therefore (3) holds. 
\end{proof}

\end{document}
